\section{Literature review}
\label{sec:literature}

\subsection{Transmission planning under the energy transition}
\label{sec:lit:planning}

Transmission is where the energy transition is most visibly constrained. The resources
that decarbonise a system sit where the wind blows and the sun shines, not where load
is concentrated, so every increment of renewable capacity implies an increment of
network --- and the network is planned centrally, built slowly, and remunerated under
rules written for a different fleet.

The modelling response has largely been optimisation. Transmission expansion
planning (TEP) treats a central planner maximising welfare subject to network
constraints, extended steadily to accommodate the transition: stochastic
formulations that make renewable hosting capacity an explicit objective
\citep{valipour2026} and multi-stage investment games for offshore grids
\citep{tao2025}. Joint generation and transmission formulations go further, and the
coordination gain is substantial: \citet{herding2024} find that coordinating the two
reduces total investment by around 10\,\% and network expansion by up to 30\,\%.
Reviews of the field \citep{gacitua2018,babatunde2020} document both its maturity
and its central assumption --- a single rational agent optimising a well-posed
objective.

That assumption is the difficulty here. In a liberalised market the generation investor
and the network planner are different agents acting on different information over
different horizons, and optimisation struggles to represent the structural and
behavioural change policy induces \citep{dyner2001}. \citet{delrio2026} make the point
concretely for Latin America: the cost of capital, not resource quality, determines
which projects are financed.

\subsection{Infrastructure delays and reliability}
\label{sec:lit:delays}

A transmission project that is planned but late is, for the duration of the delay,
indistinguishable from one that was never planned. \citet{ritter2019} simulate
delayed interconnector expansion in a high-renewable European system and show the
delay does not merely postpone benefits: it raises emissions and generation costs,
because surplus renewable energy that cannot be moved is replaced by conventional
generation elsewhere, altering the final mix. \citet{herrera2019} reach a
structurally similar conclusion for northeast Brazil, where transmission constraints
bound wind expansion directly. The mechanism is the same in both: generation
investment responds to a connection opportunity the network has not yet delivered.

The causes of delay are largely non-technical. Right-of-way disputes, licensing and
local opposition recur across jurisdictions; \citet{sun2026} quantify the last of
these, and \citet{boateng2026} generalise the point, identifying
institutional-quality thresholds below which renewable deployment responds very
differently to identical incentives. For a system such as Peru's, lead-time risk is
therefore better modelled as an institutional variable than an engineering one.

\subsection{Renewable integration in emerging economies}
\label{sec:lit:emerging}

Emerging economies face the integration problem with more expensive capital, thinner
institutions, and grids that are radial or longitudinal rather than meshed.
\citet{dasilva2025} find the regional efficiency of Brazilian renewable expansion
varying sharply by region, consistent with network access rather than resource quality
being the differentiator, while \citet{antoniolli2022}, \citet{wu2025} and
\citet{sheng2026} document rooftop and community-scale generation as mechanisms that
reduce the inter-zonal transfers a longitudinal grid must support.

The alternative to building network is using it harder. Grid-enhancing technologies
raise transfer capability without new right-of-way --- \citet{sadeghian2026} formalise
dynamic thermal line rating under chance constraints --- though remuneration
frameworks for them are generally absent. Storage plays the complementary role:
\citet{sousa2024} quantify the substitution between storage and other flexibility
resources, making the structural point that the transmission a system needs is not
independent of how much flexibility it builds.

\subsection{System dynamics in energy planning}
\label{sec:lit:sd}

System dynamics takes the opposite stance: rather than solving for the best
configuration, it simulates the configuration the system's own feedback, accumulation
and delay will produce \citep{forrester1997,sterman2000}. For electricity this is
productive because the pathologies of interest --- investment cycles, capacity gaps,
price spikes following construction lags --- are consequences of structure rather than
of optimisation failure.

The lineage is long: \citet{ford2001} explained the California capacity crisis as a
construction-delay phenomenon, \citet{arango2007} simulated alternative regulation in
Colombia, \citet{assili2008} examined capacity payments with construction delays
explicit, and \citet{simshauser2010} connected resource adequacy to investment
uncertainty. Recent work extends this to transition pathways
\citep{zapata2018,morcillo2018,luo2024}, carbon border adjustment \citep{lin2026} and
PV adoption policy \citep{oyarzun2026}.

Model confidence in system dynamics is established by a documented battery of tests
rather than a goodness-of-fit statistic \citep{barlas1996,oliva2004}, and recent
methodological work reconsiders what confidence means for models used exploratorily
rather than predictively \citep{richardson2024}. We follow that protocol explicitly
and report it in full in the supplementary material.

Where this literature is thinnest is uncertainty quantification. Scenario analysis is
near-universal; designed sampling of the input space is rare. The tools are standard
outside system dynamics --- Latin hypercube sampling \citep{mckay1992} and
regression- or variance-based global sensitivity measures \citep{homma1996} --- and
are used in power system work, for instance in \citeauthor{alvarez2026}'s
\citeyearpar{alvarez2026} probabilistic assessment of the Spanish 2030 system.
Combining them with a feedback-rich simulation model is the step this paper takes.

\subsection{Peruvian and Latin American studies, and the gap}
\label{sec:lit:gap}

Peru's restructuring is studied as a reference case: a reform that delivered
competition and improved security of supply while accumulating a high reserve margin
and fossil dependence \citep{rudnick2019,perezreyes2009}. Recent work addresses the
transition itself \citep{fiestas2024,delacruz2024} and the network
\citep{valdivia2024}. Regionally, \citet{sauma2006} established the
proactive-versus-reactive distinction in transmission valuation, while analyses of
Andean interconnection identify regulatory sovereignty rather than transfer capability
as the barrier \citep{guliev2022,sempertegui2017}.

Three gaps follow. First, the delay literature establishes that lead times matter but
not how much, relative to other uncertainties, in a common design. Second, the system
dynamics energy literature rarely quantifies uncertainty with a designed experiment,
so a small number of named scenarios cannot separate a structural result from an
artefact of one input assumption. Third, Peru is under-represented: existing studies
are either optimisation-based assessments of a renewable target
\citep{fiestas2024,delacruz2024} or operational analyses of the network
\citep{valdivia2024}, and none represents the investment-and-delay feedback between
generation and transmission that a longitudinal, gas-dependent system with suspended
auctions would be expected to exhibit. This paper addresses all three.
